\newpage
3. El uso de una lente en cierta situación se describe por la siguiente ecuación:  

\[
\frac{1}{p} + \frac{1}{-5.4p} = \frac{1}{6 \text{ cm}}
\]

Determine  

a) la distancia del objeto.  \\
b) la distancia de la imagen.  \\
c) Realice un diagrama de rayos para obtener una descripción de la imagen.  \\

Comenzamos resolviendo la ecuación dada:

\[
\frac{1}{p} + \frac{1}{-5.4p} = \frac{1}{6}.
\]


\[
\frac{1}{p} - \frac{1}{5.4p} = \frac{1}{6}.
\]



\[
\frac{1}{p} \left(1 - \frac{1}{5.4} \right) = \frac{1}{6}.
\]

Calculamos el término dentro del paréntesis:

\[
1 - \frac{1}{5.4} = \frac{5.4}{5.4} - \frac{1}{5.4} = \frac{4.4}{5.4}.
\]

Por lo tanto

\[
\frac{1}{p} \cdot \frac{4.4}{5.4} = \frac{1}{6}.
\]

\[
\frac{1}{p} = \frac{1}{6} \cdot \frac{5.4}{4.4} = \frac{5.4}{26.4}.
\]



\[
p = \frac{26.4}{5.4} \approx 4.89 \text{ cm}.
\]


La ecuación de las lentes delgadas es:

\[
\frac{1}{f} = \frac{1}{p} + \frac{1}{s_i}.
\]

Dado que ya tenemos \( p \) y sabemos que \( f = 6 \) cm:

\[
\frac{1}{6} = \frac{1}{4.89} + \frac{1}{s_i}.
\]

Simplifico:

\[
\frac{1}{s_i} = 0.1667 - 0.2044 = -0.0377.
\]

\[
s_i \approx -26.5 \text{ cm}.
\]

1. ¿Es real o virtual?
   Dado que \( s_i \) es negativo, la imagen se encuentra en el mismo lado que el objeto. Esto indica que la imagen es virtual.

2. ¿Está invertida o derecha?  
   Las imágenes virtuales en lentes convergentes siempre son derechas, aunque es más claro de ver con la respuesta del inciso M donde la magnificación es positiva $\Rightarrow$ el cociente de alturas es positivo.

3. ¿Qué tan grande es respecto al objeto?  
   La magnificación \( M \) se encuentra con:

   \[
   M = \frac{-s_i}{p} = \frac{-(-26.5)}{4.89} = \frac{26.5}{4.89} \approx 5.42.
   \]

   Como \( M > 1 \), la imagen es más grande que el objeto por un factor de 5.42 veces.

   \begin{figure}[h]
       \centering
       \includegraphics[width=1\linewidth]{traz2.jpg}
   \end{figure}